%! Algebra 2 — Unit 1, Lesson 1 — Exit Ticket (Key)
\documentclass[10pt]{article}
\usepackage{algebra2-article}
\usepackage{algebra2-boxes}
\usepackage{algebra2-key}

%=============================================================================
\begin{document}

\pageheader{Unit 1, Lesson 1}{Exit Ticket --- Key}
\namedateperiod

\vspace{0.3in}

\begin{enumerate}[leftmargin=*, itemsep=1.8em]

    % ── Question 1 ──────────────────────────────────────────────────────────
    \item Evaluate $2x^2 + 3x - 4$ when $x = -3$. \textbf{Show your work.}

    \vspace{0.1in}
    \ansline{$2(-3)^2 + 3(-3) - 4$}
    \ansline{$= 2(9) - 9 - 4$}
    \ansline{$= 18 - 9 - 4$}
    \ansline{$= 5$}

    \vspace{0.05in}
    Answer: \ans{$5$}

    \begin{teachernote}
        Watch for students who write $2(-3)^2 = 2(-9) = -18$ (failing to square before applying the negative, or squaring the negative without the coefficient). Also watch for sign errors on $3(-3) = 9$ instead of $-9$.
    \end{teachernote}

    % ── Question 2 ──────────────────────────────────────────────────────────
    \item Simplify $\sqrt{48}$. \textbf{Show your work.}

    \vspace{0.1in}
    \ansline{$\sqrt{48} = \sqrt{16 \cdot 3}$}
    \ansline{$= \sqrt{16}\cdot\sqrt{3}$}
    \ansline{$= 4\sqrt{3}$}

    \vspace{0.05in}
    Answer: \ans{$4\sqrt{3}$}

    \begin{teachernote}
        Common error: $\sqrt{48} = \sqrt{4 \cdot 12} = 2\sqrt{12}$ (not fully simplified — students must find the \emph{largest} perfect-square factor). Push students to verify by asking whether $12$ contains another perfect-square factor.
    \end{teachernote}

    % ── Question 3 (SOL-style MC) ────────────────────────────────────────────
    \item Which expression is equivalent to $\displaystyle\frac{x^5 \cdot x^2}{x^4}$?

    \vspace{0.25in}
    \begin{tabularx}{\linewidth}{Y Y Y Y}
        \textbf{A.}\quad $x^{\frac{3}{4}}$
        & \ans{\textbf{B.}\quad $x^{3}$}
        & \textbf{C.}\quad $x^{6}$
        & \textbf{D.}\quad $x^{11}$
    \end{tabularx}

    \vspace{0.2in}
    \ans{$\dfrac{x^5 \cdot x^2}{x^4} = \dfrac{x^{5+2}}{x^4} = \dfrac{x^7}{x^4} = x^{7-4} = x^3$} \quad $\Rightarrow$ \ans{\textbf{B}}

    \begin{teachernote}
        \textbf{Distractor analysis:} \textbf{A} — divided the final exponent by $4$ instead of subtracting; \textbf{C} — multiplied instead of adding in the numerator ($5 \times 2 = 10$, then $10-4=6$); \textbf{D} — added all three exponents ($5+2+4=11$).
    \end{teachernote}

\end{enumerate}

\end{document}
